\begin{table*}[htbp]
\centering
\scriptsize
\renewcommand{\arraystretch}{1.35}
\begin{NiceTabularX}{\textwidth}{
>{\centering\arraybackslash}m{0.9cm}
>{\centering\arraybackslash}m{1.8cm}
>{\centering\arraybackslash}m{2.1cm}
X X X X
}[hvlines]
核心维度 & 细分落地 & 核心主体 & 代表性成果 & 技术方案 & 核心价值 & 不足 \\
\Block{3-1}{学术侧} &
\Block{3-1}{机密计算、TEE 与云端 AI 隔离边界验证} &
\Block{3-1}{ETH Zurich、TU Berlin、CISPA、KU Leuven、Cambridge、Imperial、EPFL 等系统安全团队} &
SEVered 证明恶意 hypervisor 可从 AMD SEV 加密 VM 中抽取明文内存 \cite{severed2018}。 &
构造恶意页映射和服务交互，把加密 VM 的内存内容转化为可抽取明文。 &
说明 CVM/TEE 是 AI 云底座的重要机制，但不能视为一次性可信。 &
多数攻击针对 CPU TEE，和 GPU/NPU 机密执行的直接结合仍需扩展。 \\
& & &
One Glitch to Rule Them All 通过 AMD Secure Processor 故障注入破坏 SEV 信任根，可解密 VM 内存并伪造证明相关材料 \cite{oneglitch2021}。 &
对安全协处理器进行电压故障注入，绕过固件信任根和密钥保护。 &
推动云端 AI 证明链从“CPU 支持”扩展到固件、异常处理、设备路由和平台初始化。 &
防护依赖芯片商、OEM、云厂商同步更新固件和证明策略。 \\
& & &
WeSee 利用 SEV-SNP 的 \#VC 异常处理路径，展示恶意 hypervisor 仍可影响机密 VM 执行完整性 \cite{wesee2024}。 &
注入恶意虚拟化异常，引导 guest 内核执行非预期数据复制、状态破坏或代码路径。 &
为高风险 Agent 平台提供“持续攻击-验证-修补”的安全论证方式。 &
真实云平台内部配置不完全公开，第三方可复验性有限。 \\
\Block{3-1}{学术侧} &
\Block{3-1}{SEV-SNP、远程证明与平台初始化细节安全} &
\Block{3-1}{ETH Zurich、CISPA、欧洲形式化方法与云安全团队；AMD PSIRT 协调披露生态} &
AMD-SB-3034 披露 SEV-SNP Routing Misconfiguration，MMIO 路由和 AMD Secure Processor 固件锁定检查缺陷可影响 guest 完整性 \cite{amdsb3034}。 &
检查平台初始化、RMP、MMIO 路由、ASP/PSP 固件锁定和 TCB 版本传播。 &
把“是否运行在 SEV-SNP 中”细化为“以何种固件、路由、证明和配置运行”。 &
证明链过长，普通部署方难以理解每个 TCB 字段和固件版本含义。 \\
& & &
SEV-SNP 软件接口形式化分析覆盖远程证明、密钥派生、页面交换和迁移，并指出平台无关消息设计带来的证明完整性弱点 \cite{formalsevsnp2024}。 &
用符号模型刻画证明、迁移、密钥派生协议，验证 secrecy、authentication、freshness 等性质。 &
支撑公共部门和高风险行业在采购 CVM/机密 AI 服务时提出可审计要求。 &
形式化模型通常抽象硬件实现细节，难完全覆盖实际固件和主板配置。 \\
& & &
SEVered/WeSee 等连续工作共同构成欧洲对 SEV 系列机密计算边界的持续回归测试 \cite{severed2018,wesee2024}。 &
将漏洞公告、固件 TCB 值和云端证明策略纳入同一条安全链。 &
强化 Agent 平台对底层证明材料、固件版本和异常路径的证据化治理。 &
漏洞修复依赖 BIOS/固件发布节奏，云租户缺乏直接控制权。 \\
\Block{3-1}{学术侧} &
\Block{3-1}{工业边缘可信连接、设备证明与数据空间安全} &
\Block{3-1}{Fraunhofer、Industrial Data Space/International Data Spaces 生态、欧洲嵌入式/IoT 安全团队} &
Industrial Data Space Trusted Connector 将可信边缘网关用于跨组织工业数据交换，强调容器隔离、身份、策略执行和远程证明 \cite{trustedconnector2018}。 &
在边缘网关中结合 TPM/TEE、容器、策略引擎和远程证明。 &
更贴合工业 Agent 接触传感器、PLC、网关、运维账户和现场数据的部署形态。 &
工业现场设备生命周期长，旧设备缺少 TEE/TPM/安全元件。 \\
& & &
PSA Certified/Arm TrustZone 生态把 Root of Trust、证明 API、加密 API 和安全存储抽象为 IoT/边缘设备可复用框架 \cite{psacertified2026,armplatformsecurity2026}。 &
以 PSA RoT、TrustZone、Secure Boot、Attestation、Secure Storage 为设备侧信任基元。 &
支撑“数据不出域、节点可证明、策略可审计”的欧洲工业数据空间目标。 &
跨厂商网关、PLC、工业协议和数据空间策略互操作复杂。 \\
& & &
TrustZone 运行时完整性测量、TEE 安全缓存等研究探索在边缘设备上保护敏感数据和本地 AI/TinyML 工作负载 \cite{kevlartz2021,pdrima2025}。 &
对 TEE 内核、TA、安全调用和本地缓存进行运行时测量与完整性审计。 &
为边缘 AI、TinyML 和工业 Agent 的设备身份与固件可信提供基础。 &
边缘证明仍难覆盖物理篡改、现场维护绕过和长期补丁缺口。 \\
\Block{3-1}{产业侧} &
\Block{3-1}{主权算力、合规 AI 与公共超算基础设施} &
\Block{3-1}{European Commission、EuroHPC JU、成员国超算中心、AI Factories、JUPITER/LUMI/BSC 等生态} &
European Chips Act 强调增强欧洲半导体生态、供应链安全、韧性、数字主权和降低第三方依赖 \cite{euchipsact2023}。 &
用法规和公共投资把芯片供应、AI 算力、数据治理和高风险系统合规连接起来。 &
让 AI 硬件底座成为公共部门和产业主权能力，而不只是商业云资源。 &
高端 GPU/加速器、云软件栈和超大模型平台仍较依赖外部生态。 \\
& & &
EU AI Act 以风险分级方式约束高风险 AI 系统，要求风险管理、数据治理、技术文档、记录、人类监督、准确性、鲁棒性和网络安全 \cite{euaiact2024}。 &
通过 AI Factories 把超算中心、大学、SME、创业公司、产业和公共部门接入统一服务网络。 &
为高风险 Agent 在医疗、金融、交通、能源、制造等场景部署提供合规接口。 &
合规要求清晰但执行成本高，SME 和科研团队负担较重。 \\
& & &
EuroHPC AI Factories 已形成 19 个 AI Factories 和 13 个 Antennas，为欧洲用户提供 AI 优化超算和支持服务 \cite{eurohpcAIFactories2026,ecAIFactories2026}。 &
在采购、访问控制、支持服务和合规文档中嵌入可信 AI 要求。 &
以公共采购和监管要求反向推动安全证明、审计和长期维护能力。 &
法规语言到芯片级安全指标之间仍需标准和测评工具衔接。 \\
\Block{3-1}{产业侧} &
\Block{3-1}{工业、汽车与嵌入式安全芯片底座} &
\Block{3-1}{Infineon、STMicroelectronics、NXP、Bosch、Siemens、Arm 欧洲生态} &
Infineon OPTIGA Trust 面向 MCU/MPU/应用处理器提供安全认证、安全通信、安全更新和访问管理 \cite{infineonoptiga2026}。 &
用安全元件、TPM/HSM、唯一设备身份和预置证书绑定工业设备。 &
适合工业 Agent 进入现场系统时建立设备身份、固件可信和运维权限边界。 &
产品级安全能力分散在不同芯片、工具链和证书体系中。 \\
& & &
STSAFE-A110 作为安全元件提供认证、安全数据管理、TLS 握手、签名验证、安全启动/固件升级、抗故障和抗侧信道能力 \cite{stsafea1102026}。 &
通过 Secure Boot、固件签名、TLS、密钥隔离、计数器和安全更新保护边缘节点。 &
保护 PLC、边缘网关、传感器、机器人和汽车电子中的模型、配置和现场数据。 &
端侧算力、内存和功耗限制使复杂 Agent 难完全本地运行。 \\
& & &
PSA Certified/Arm 平台安全生态为 IoT/边缘芯片建立多级安全评估、RoT、证明、加密和存储 API \cite{psacertified2026,armplatformsecurity2026}。 &
以 PSA/IEC 62443/CRA 等框架把芯片安全能力转译成认证和采购语言。 &
与欧洲功能安全、网络安全和长期维护传统形成协同。 &
硬件安全功能需要系统集成方正确配置，否则容易停留在“可用但未启用”。 \\
\Block{3-1}{产业侧} &
\Block{3-1}{工业现场安全运营、合规采购与长期维护} &
\Block{3-1}{Siemens、Bosch Rexroth、Schneider Electric、工业软件厂商、ProductCERT/PSIRT 生态} &
Siemens 工业网络安全方案强调从管理系统到现场层、从访问控制到数据流的防护，并采用符合 IEC 62443 的 defense-in-depth 思路 \cite{siemensindustrialcyber2026}。 &
将 Agent 可触达的工业资产分层：现场设备、控制网络、边缘网关、SCADA、MES/ERP 和云服务。 &
使 Agent 安全从模型越狱治理扩展到工业现场可运行、可审计、可维护。 &
工业现场异构且生命周期长，补丁窗口和停机成本限制修复速度。 \\
& & &
Siemens 将 plant security、network security、system integrity 与 Zero Trust 原则结合，用于保护控制器、HMI、SCADA 和远程维护接口 \cite{siemensindustrialcyber2026}。 &
结合资产清单、访问控制、远程维护审计、系统完整性和安全补丁流程。 &
让硬件安全、工业协议安全和法规合规进入同一套验收语言。 &
供应商安全承诺、集成商配置和业主运维能力之间可能存在断点。 \\
& & &
工业厂商通过 ProductCERT、固件更新、工业安全公告和认证材料支撑 NIS2、CRA、IEC 62443 等合规要求 \cite{siemensindustrialcyber2026,infineonoptiga2026}。 &
在采购规范中要求安全启动、可更新固件、漏洞披露、生命周期支持和证明材料。 &
有利于能源、交通、制造和公共设施中的高风险 Agent 形成责任链。 &
工业 Agent 的自治决策、跨系统联动和责任归属仍缺少成熟认证范式。 \\
\end{NiceTabularX}
\caption{欧洲 AI 硬件与芯片安全研究现状：以方向为主轴、主体为支撑的结构化总结}
\label{tab:europe-ai-hardware-security-nicematrix}
\end{table*}
